\chapter{Discrete-Time Derivative Pricing}

\begin{objectives}
Price derivatives by replication in one- and multi-period trees, understand
completeness and early exercise, and derive convergence toward continuous-time
models.
\end{objectives}

\section{One-period binomial economy}

Let a stock cost $S_0$ and move to $S_u=uS_0$ or $S_d=dS_0$, with
$d<R<u$ where $R$ is the risk-free gross return. A claim pays $H_u,H_d$. Choose
$\Delta$ shares and cash $B$ so
\[
\Delta S_u+RB=H_u,\qquad
\Delta S_d+RB=H_d.
\]
Solving,
\[
\Delta=\frac{H_u-H_d}{S_0(u-d)},\qquad
B=\frac{uH_d-dH_u}{R(u-d)}.
\]
The no-arbitrage value is $V_0=\Delta S_0+B$.

Define risk-neutral probability
\[
q=\frac{R-d}{u-d}\in(0,1).
\]
Then
\[
V_0=\frac1R[qH_u+(1-q)H_d].
\]
The physical up probability never entered replication. It matters for risk and
forecasting, not the unique price in this complete tree.

\begin{theorem}[No-arbitrage condition]
The one-period stock--bond binomial market is arbitrage free if and only if
$d<R<u$. It is complete when $u\neq d$.
\end{theorem}
\begin{proof}
If $R\leq d$, borrowing to buy stock has nonnegative payoff in both states and a
strict gain in at least one; similarly $R\geq u$ supports the opposite trade. If
$d<R<u$, $q\in(0,1)$ gives strictly positive state prices, excluding arbitrage.
The two linearly independent securities span the two states.
\end{proof}

\section{Multi-period trees}

At each node, value satisfies backward recursion
\[
V_{t}=\frac1R\left(qV_{t+1}^{u}+(1-q)V_{t+1}^{d}\right).
\]
The replicating delta is recomputed at every node, making the hedge dynamic and
self-financing. In a recombining $n$-step tree a European payoff has
\[
V_0=R^{-n}\sum_{j=0}^n
\binom{n}{j}q^j(1-q)^{n-j}
H(S_0u^jd^{n-j}).
\]

The Cox--Ross--Rubinstein parameterization for step $\Delta t$ uses
$u=\e^{\sigma\sqrt{\Delta t}}$, $d=u^{-1}$, and
$R=\e^{(r-q_d)\Delta t}$. As the grid is refined, the log-price random walk
converges weakly to geometric Brownian motion under appropriate conditions.

\section{American exercise}

At an American-option node,
\[
V_t=\max\left\{h(S_t),
\frac1R\E_t^\Q[V_{t+1}]\right\}.
\]
This is an optimal stopping problem. The value process is the smallest
supermartingale dominating the discounted exercise payoff, the Snell envelope.

\begin{proposition}[No early exercise for a non-dividend call]
With nonnegative rates and no dividends, it is never optimal to exercise an American
call before maturity.
\end{proposition}
\begin{proof}
By put--call parity and $P_t\geq0$,
$C_t\geq S_t-KP(t,T)\geq S_t-K$. Continuation value is at least intrinsic value, and
exercise destroys remaining optionality and the financing benefit of delaying $K$.
\end{proof}

American puts and dividend-paying calls can exercise early. A tree must align
dividend dates, exercise opportunities, and contract conventions.

\section{Incomplete trees and bounds}

If there are more states than independent traded securities, replication may fail.
Sub- and super-replication define no-arbitrage price bounds:
\[
\pi_{\mathrm{sub}}(H)\leq V_0\leq\pi_{\mathrm{super}}(H).
\]
A utility indifference price, minimal entropy measure, or calibrated model selects a
point only by adding criteria beyond no arbitrage.

\section{Problems}

\begin{exercise}[Core: two-step call]
Price and replicate a two-step European call. Record delta and cash at every node
and verify terminal payoffs.
\end{exercise}
\begin{exercise}[Core: American put]
Construct a two-step example where early put exercise is optimal and explain the
roles of rate and moneyness.
\end{exercise}
\begin{exercise}[Advanced: convergence]
Match the first two conditional moments of the CRR log return and show their limits
as $\Delta t\downarrow0$.
\end{exercise}


\chapter{Black--Scholes--Merton Theory and Greeks}

\begin{objectives}
Derive the pricing PDE and closed form, understand hedging and Greeks, and identify
exactly which conclusions rely on which assumptions.
\end{objectives}

\section{Market assumptions}

In the benchmark model,
\[
\dd S_t=(\mu-q)S_t\dd t+\sigma S_t\dd W_t^\Prob,\qquad
\dd B_t=rB_t\dd t,
\]
with constant $\sigma>0$, deterministic $r,q$, continuous trading, no jumps,
frictionless markets, unrestricted admissible trading, and a European payoff. The
stock including dividends earns drift $\mu$.

Under the pricing measure,
\[
\dd S_t=(r-q)S_t\dd t+\sigma S_t\dd W_t^\Q.
\]

\section{Delta-hedging derivation}

For value $V(t,S)$, Itô's formula gives
\[
\dd V=\left(V_t+(\mu-q)SV_S+\frac12\sigma^2S^2V_{SS}\right)\dd t
+\sigma SV_S\dd W.
\]
Hold one derivative and short $\Delta=V_S$ shares. After incorporating dividend on
the stock hedge, stochastic risk cancels. No arbitrage makes the locally riskless
portfolio earn $r$, yielding
\[
V_t+(r-q)SV_S+\frac12\sigma^2S^2V_{SS}-rV=0.
\]
The physical drift $\mu$ cancels because the source of diffusion risk is spanned.

\section{Risk-neutral valuation}

Feynman--Kac gives
\[
V(t,S)=\e^{-r\tau}\E^\Q[h(S_T)\mid S_t=S],
\qquad \tau=T-t,
\]
with
\[
\log S_T\mid S_t
\sim N\left(\log S+(r-q-\tfrac12\sigma^2)\tau,\sigma^2\tau\right).
\]

\begin{theorem}[European call and put]
\[
C=S\e^{-q\tau}\Phi(d_1)-K\e^{-r\tau}\Phi(d_2),
\]
\[
P=K\e^{-r\tau}\Phi(-d_2)-S\e^{-q\tau}\Phi(-d_1),
\]
where
\[
d_1=\frac{\log(S/K)+(r-q+\frac12\sigma^2)\tau}
{\sigma\sqrt\tau},\qquad
d_2=d_1-\sigma\sqrt\tau.
\]
\end{theorem}
\begin{proof}
Write the call as
$\e^{-r\tau}(\E^\Q[S_T\ind_{S_T>K}]
-K\Prob^\Q(S_T>K))$. Standardizing the lognormal threshold gives
$\Prob^\Q(S_T>K)=\Phi(d_2)$. Exponential tilting of the normal density gives
$\E^\Q[S_T\ind_{S_T>K}]=S\e^{(r-q)\tau}\Phi(d_1)$. The put follows from parity.
\end{proof}

\section{Greeks}

For a call,
\[
\Delta_C=\e^{-q\tau}\Phi(d_1),\qquad
\Gamma=\frac{\e^{-q\tau}\phi(d_1)}{S\sigma\sqrt\tau},
\]
\[
\mathcal V=\frac{\partial C}{\partial\sigma}
=S\e^{-q\tau}\phi(d_1)\sqrt\tau,
\]
\[
\Theta_C=
-\frac{S\e^{-q\tau}\phi(d_1)\sigma}{2\sqrt\tau}
-rK\e^{-r\tau}\Phi(d_2)
+qS\e^{-q\tau}\Phi(d_1),
\]
\[
\rho_C=K\tau\e^{-r\tau}\Phi(d_2).
\]
Greek units matter: vega above is per absolute volatility unit, so divide by $100$
for a one-volatility-point move.

The PDE implies a theta--gamma--carry relation. For a delta-hedged option, time decay
compensates for expected gamma gains under the model. Large gamma near expiry makes
discrete hedging difficult exactly when theta is largest.

\section{Discrete hedging error}

With true realized quadratic variation different from model variance, a continuous
delta-hedge approximation gives
\[
\Pi_T-\Pi_0\e^{rT}
\approx\frac12\int_0^T
\e^{r(T-t)}\Gamma_tS_t^2
(\dd\langle\log S\rangle_t-\sigma_{\mathrm{model}}^2\dd t)
\]
before costs and jumps. Discrete rebalancing adds error of order depending on mesh,
payoff regularity, and strategy. Transaction costs grow with trading frequency, so
continuous hedging is not an operational limit without a cost model.

\section{Digital options and density}

A cash-or-nothing call pays $\ind_{\{S_T>K\}}$ and is worth
$\e^{-r\tau}\Phi(d_2)$. Under regularity,
\[
-\frac{\partial C(K,T)}{\partial K}
=P(t,T)\Q(S_T>K),\qquad
\frac{\partial^2 C(K,T)}{\partial K^2}
=P(t,T)f^\Q_{S_T}(K).
\]
Thus a continuum of option prices reveals the risk-neutral terminal distribution.
Noise and discrete strikes make differentiation an ill-posed estimation problem
requiring arbitrage-aware smoothing.

\section{Problems}

\begin{exercise}[Core: Greeks]
Differentiate the call formula to obtain delta, gamma, and vega. Use the identity
$S\e^{-q\tau}\phi(d_1)=K\e^{-r\tau}\phi(d_2)$.
\end{exercise}
\begin{exercise}[Core: PDE]
Verify that the call formula satisfies the PDE, terminal condition, and boundary
behavior.
\end{exercise}
\begin{exercise}[Advanced: digital replication]
Show that a tight call spread converges to a digital payoff and quantify the finite
strike-spacing approximation.
\end{exercise}


\chapter{Volatility Surfaces, Local and Stochastic Volatility}

\begin{objectives}
Construct arbitrage-consistent volatility surfaces, derive local volatility, study
stochastic-volatility dynamics, and distinguish calibration fit from hedging validity.
\end{objectives}

\section{Why volatility is a surface}

Market implied volatility depends on strike and maturity. Equity index smiles often
show expensive downside protection; term structure changes with event and regime.
Quoting implied volatility normalizes option prices but preserves model dependence.

Useful coordinates include log forward moneyness
$k=\log(K/F_{t,T})$, delta, and total variance
$w(k,T)=\sigma_{\mathrm{imp}}^2(k,T)T$. Interpolating raw volatility independently
can create arbitrage even when input quotes are valid.

\section{Static-arbitrage restrictions}

For fixed $T$, undiscounted call price is decreasing and convex in $K$:
\[
\partial_K C\leq0,\qquad \partial_{KK}C\geq0.
\]
It lies between intrinsic lower bounds and the prepaid forward. Across maturities,
calendar consistency is naturally expressed for option prices under aligned forwards
and discounting; total variance should not decrease under simplified settings.

Butterfly arbitrage corresponds to negative risk-neutral density. Calendar arbitrage
corresponds to a more valuable shorter-dated claim under conditions where the
longer-dated claim contains it.

\section{Dupire local volatility}

Assume under $\Q$
\[
\frac{\dd S_t}{S_t}=(r_t-q_t)\dd t+\sigma_{\mathrm{loc}}(t,S_t)\dd W_t.
\]
Given a sufficiently smooth arbitrage-free call surface,
\[
\sigma_{\mathrm{loc}}^2(T,K)=
\frac{\partial_T C+(r-q)K\partial_K C+qC}
{\frac12K^2\partial_{KK}C}
\]
for deterministic rates and yields under the stated convention.

Local volatility exactly matches the one-time marginal distributions encoded by the
surface. It does not necessarily reproduce joint dynamics, forward smiles, or
realistic spot--volatility co-movement. Numerical differentiation amplifies quote
noise, especially through $\partial_{KK}C$.

\section{Heston stochastic volatility}

The Heston model under $\Q$ is
\[
\frac{\dd S_t}{S_t}=(r-q)\dd t+\sqrt{v_t}\dd W_t^S,
\]
\[
\dd v_t=\kappa(\theta-v_t)\dd t+\xi\sqrt{v_t}\dd W_t^v,\qquad
\dd\langle W^S,W^v\rangle_t=\rho\dd t.
\]
Mean reversion is controlled by $\kappa$, long-run variance by $\theta$, volatility
of variance by $\xi$, and skew by $\rho$. The Feller condition
$2\kappa\theta\geq\xi^2$ keeps zero unattainable under one common classification but
is not required for every numerical scheme or admissible calibration.

The characteristic function is affine and supports Fourier pricing. Parameters can
be poorly identified from a narrow option panel: $\kappa$ and $\theta$, for example,
may trade off while producing similar maturities.

\section{SABR}

For forward $F$,
\[
\dd F_t=\alpha_tF_t^\beta\dd W_t^1,\qquad
\dd\alpha_t=\nu\alpha_t\dd W_t^2,\qquad
\dd\langle W^1,W^2\rangle_t=\rho\dd t.
\]
SABR is widely used to interpolate smiles, especially in rates. $\beta$ controls
backbone, $\nu$ smile curvature, and $\rho$ skew. Common asymptotic implied-volatility
formulas can be inaccurate at long maturity, extreme strikes, or near boundaries.

\section{Jumps and Lévy models}

Merton's jump diffusion adds compound Poisson lognormal jumps. Variance gamma and
CGMY use infinite-activity Lévy processes. Jumps create short-maturity skew and
kurtosis that continuous diffusions struggle to match.

Risk-neutral jump intensity and distribution are not physical jump forecasts.
Hedging jump risk is incomplete when no traded asset spans it; calibration chooses
prices consistent with selected instruments rather than manufacturing completeness.

\section{Calibration}

For quotes $P_i^{\mathrm{mkt}}$ and weights $w_i$,
\[
\widehat\theta=
\argmin_{\theta\in\Theta}
\sum_iw_i(P_i^{\mathrm{model}}(\theta)-P_i^{\mathrm{mkt}})^2
+\lambda\mathcal R(\theta).
\]
Price-error weighting, vega weighting, relative errors, and bid--ask normalization
answer different questions. A robust workflow checks:

\begin{itemize}
\item no-arbitrage cleaning and quote timestamps;
\item multiple starting values and parameter bounds;
\item residuals by strike and maturity;
\item stability across days;
\item prices of withheld instruments;
\item hedge performance and scenario behavior;
\item parameter uncertainty and numerical tolerance.
\end{itemize}

\begin{warning}
Six decimal places in a calibrated parameter are usually optimizer output, not six
decimal places of economic information.
\end{warning}

\section{Problems}

\begin{exercise}[Core: convexity]
Use a butterfly portfolio to prove convexity of call price in strike.
\end{exercise}
\begin{exercise}[Core: Heston moments]
Derive $\E[v_t\mid v_0]$ and the long-run mean. Discuss what changes under the
physical measure.
\end{exercise}
\begin{exercise}[Advanced: Dupire]
Derive the Dupire equation from the forward equation for the transition density and
integration by parts.
\end{exercise}


\chapter{Exotic Derivatives and Incomplete-Market Hedging}

\begin{objectives}
Analyze path-dependent payoffs, apply change-of-numeraire methods, value options with
exercise features, and measure hedging risk in incomplete markets.
\end{objectives}

\section{Change of numeraire}

Let $N_t>0$ be a traded numeraire. There exists an associated measure $\Q^N$ under
which every traded price divided by $N$ is a martingale:
\[
\frac{V_t}{N_t}=
\E_t^{\Q^N}\left[\frac{V_T}{N_T}\right].
\]
Choosing a zero-coupon bond as numeraire gives a forward measure, under which the
corresponding forward price is a martingale. This simplifies caplet, swaption, and
quanto calculations by matching measure to payoff.

\section{Asian options}

An arithmetic-average call pays
\[
\left(\frac1n\sum_{j=1}^nS_{t_j}-K\right)^+.
\]
The arithmetic average of lognormal variables is not lognormal, so no Black--Scholes
closed form exists. Geometric-average options are tractable and provide control
variates. Bounds follow from Jensen and comonotonic constructions.

\section{Barrier options}

A down-and-out call ceases to exist if spot reaches barrier $H<S_0$. Continuous
monitoring differs from daily monitoring because crossings can occur between
observations. Reflection-principle formulas exist in Black--Scholes for standard
contracts. Greeks become sharp near the barrier and discrete hedge error grows.

In--out parity gives
\[
V_{\mathrm{vanilla}}=V_{\mathrm{knock\mbox{-}in}}+
V_{\mathrm{knock\mbox{-}out}}
\]
when rebates and monitoring conventions align.

\section{Lookbacks and cliquets}

Lookback payoffs depend on running maximum or minimum. Cliquets lock periodic
returns subject to local and global caps or floors. Their values depend strongly on
the joint forward volatility surface, not only today's marginal smiles. A model
calibrated to vanillas can therefore disagree materially on exotics.

\section{Bermudan options and least-squares Monte Carlo}

At exercise dates $t_m$, a Bermudan value is
\[
V_{t_m}=\max\{h(X_{t_m}),\E^\Q[D_{m,m+1}V_{t_{m+1}}\mid X_{t_m}]\}.
\]
Least-squares Monte Carlo regresses discounted continuation cash flows on basis
functions of state among relevant paths, then compares fitted continuation with
exercise value.

Using the same paths to fit and value creates in-sample bias. Independent valuation
paths, cross-fitting, basis sensitivity, and dual upper bounds improve credibility.
The exercise rule should respect all state variables, not spot alone when rates or
volatility are stochastic.

\section{Variance and volatility derivatives}

A variance swap pays realized variance minus fixed strike. Under ideal continuous
sampling and diffusion,
\[
\int_0^T\sigma_t^2\dd t
=2\int_0^T\frac{\dd S_t}{S_t}
-2\log\frac{S_T}{S_0}
\]
after appropriate drift terms. A log contract can be replicated by a continuum of
out-of-the-money options, giving a model-light fair variance strike under carefully
stated jump and discretization conventions.

Volatility swap payoff uses square root of variance, so Jensen creates a convexity
adjustment; it cannot be priced by simply taking the square root of variance-swap
strike.

\section{Quanto and correlation risk}

A quanto pays an asset return in a different currency at a fixed exchange rate.
Under the domestic pricing measure, covariance between the foreign asset and FX
changes its drift. This quanto adjustment makes correlation a first-order input.
Cross-asset derivatives reveal why specifying only marginal models is insufficient.

\section{Hedging in incomplete markets}

When risk is unspanned, exact replication is impossible. A minimum-variance hedge
chooses strategy $\phi$ to minimize
\[
\E[(H-V_0-G_T(\phi))^2].
\]
The Galtchouk--Kunita--Watanabe decomposition writes payoff into initial value,
stochastic integral with traded assets, and an orthogonal residual. Utility
indifference and risk-measure hedges choose different objectives and therefore
different prices.

\section{Problems}

\begin{exercise}[Core: geometric Asian]
Derive the distribution of a discretely sampled geometric average under
Black--Scholes and write its call value.
\end{exercise}
\begin{exercise}[Core: variance replication]
Show how a twice differentiable payoff can be statically represented using cash,
forward, and a continuum of calls and puts. Apply it to $-\log S_T$.
\end{exercise}
\begin{exercise}[Advanced: LSM]
Implement least-squares Monte Carlo for an American put and report lower-bound
stability across bases, paths, and time steps.
\end{exercise}
